%%%%%%%%%%%%%%%%%%%%%%%%%%%%%%%%%%%%%%%%%%%%%%%%%%%%%%%%%%%%%%%%%%%%%%%%%%%%%%%%
%% Capitulo 4: Documento de Arquitetura de Software (SAD)
%% Sistema: PostoSmart - Automação e Gestão de Postos de Combustíveis
%% Autor: Diogo Santos Pires Jandiroba
%%%%%%%%%%%%%%%%%%%%%%%%%%%%%%%%%%%%%%%%%%%%%%%%%%%%%%%%%%%%%%%%%%%%%%%%%%%%%%%%

\chapter{Documento de Arquitetura de Software (SAD)}
\label{cap:arquitetura_software}

\section{Visão Geral e Estilo Arquitetural}

A arquitetura do \textbf{PostoSmart} adota o padrão \textit{Edge-First Event-Driven} (borda resiliente orientada a eventos). Como postos de combustíveis operam em condições agressivas com instabilidade de conectividade WAN, toda a lógica de autorização de bombas, controle de pulsos e frente de caixa executa localmente em um appliance local (\textit{Edge Server}), sincronizando assincronamente com o cluster de nuvem.

\section{Visão Lógica (Componentes do Sistema)}

A \autoref{fig:visao_logica_posto} ilustra a decomposição em subsistemas modulares com alta coesão e baixo acoplamento.

% PLACEHOLDER STARUML v7.0: Visão Lógica (Componentes)
% Ferramenta: StarUML v7.0 -> Model -> Add Diagram -> Component Diagram
% Arquivo: Imagens/arch_visao_logica.png (Exportar em PNG 300 DPI ou PDF vetorial)
% Descomente a linha abaixo apos salvar o arquivo no diretorio Imagens/:
% \incluirdiagrama{Imagens/arch_visao_logica.png}{Visão Lógica da Arquitetura Edge-First do PostoSmart}{fig:visao_logica_posto}
\begin{figure}[htbp]
    \centering
    \fbox{\parbox{0.9\textwidth}{\centering\vspace{1cm}%
        \textbf{[Visão Lógica -- Diagrama de Componentes -- StarUML v7.0]}\\%
        \textit{Modelar no StarUML v7.0 e exportar para: \texttt{Imagens/arch\_visao\_logica.png}}\\%
        \small Menu: \texttt{File -> Export Diagram as -> PNG...} (Fundo branco, alta resolução 300 DPI)\\%
        \vspace{0.3cm}%
        \footnotesize Para ativar a imagem real, descomente no arquivo \texttt{.tex}:\\
        \texttt{\textbackslash incluirdiagrama\{Imagens/arch\_visao\_logica.png\}\{Visão Lógica\}\{fig:visao_logica_posto\}}%
    \vspace{1cm}}}
    \caption{Visão Lógica da Arquitetura Edge-First do PostoSmart (Placeholder StarUML v7.0)}
    \label{fig:visao_logica_posto}
\end{figure}

\section{Visão de Processos e Concorrência}

A \autoref{fig:visao_processos_posto} detalha o processamento concorrente em tempo real dos pulsos emitidos pelas bombas de combustível.

% PLACEHOLDER STARUML v7.0: Visão de Processos e Concorrência
% Ferramenta: StarUML v7.0 -> Model -> Add Diagram -> Sequence Diagram ou Activity Diagram
% Arquivo: Imagens/arch_visao_processos.png (Exportar em PNG 300 DPI ou PDF vetorial)
% Descomente a linha abaixo apos salvar o arquivo no diretorio Imagens/:
% \incluirdiagrama{Imagens/arch_visao_processos.png}{Visão de Processos do Pipeline de Telemetria de Pista}{fig:visao_processos_posto}
\begin{figure}[htbp]
    \centering
    \fbox{\parbox{0.9\textwidth}{\centering\vspace{1cm}%
        \textbf{[Visão de Processos e Concorrência -- StarUML v7.0]}\\%
        \textit{Modelar no StarUML v7.0 e exportar para: \texttt{Imagens/arch\_visao\_processos.png}}\\%
        \small Menu: \texttt{File -> Export Diagram as -> PNG...} (Fundo branco, alta resolução 300 DPI)\\%
        \vspace{0.3cm}%
        \footnotesize Para ativar a imagem real, descomente no arquivo \texttt{.tex}:\\
        \texttt{\textbackslash incluirdiagrama\{Imagens/arch\_visao\_processos.png\}\{Visão de Processos\}\{fig:visao_processos_posto\}}%
    \vspace{1cm}}}
    \caption{Visão de Processos do Pipeline de Telemetria de Pista (Placeholder StarUML v7.0)}
    \label{fig:visao_processos_posto}
\end{figure}

\section{Visão Física e de Implantação}

A \autoref{fig:visao_implantacao_posto} descreve a infraestrutura física de pista, cabeamento industrial e ambiente corporativo em nuvem.

% PLACEHOLDER STARUML v7.0: Visão Física e de Implantação
% Ferramenta: StarUML v7.0 -> Model -> Add Diagram -> Deployment Diagram
% Arquivo: Imagens/arch_visao_implantacao.png (Exportar em PNG 300 DPI ou PDF vetorial)
% Descomente a linha abaixo apos salvar o arquivo no diretorio Imagens/:
% \incluirdiagrama{Imagens/arch_visao_implantacao.png}{Visão Física de Implantação e Conectividade Industrial}{fig:visao_implantacao_posto}
\begin{figure}[htbp]
    \centering
    \fbox{\parbox{0.9\textwidth}{\centering\vspace{1cm}%
        \textbf{[Visão Física e de Implantação -- Deployment Diagram -- StarUML v7.0]}\\%
        \textit{Modelar no StarUML v7.0 e exportar para: \texttt{Imagens/arch\_visao\_implantacao.png}}\\%
        \small Menu: \texttt{File -> Export Diagram as -> PNG...} (Fundo branco, alta resolução 300 DPI)\\%
        \vspace{0.3cm}%
        \footnotesize Para ativar a imagem real, descomente no arquivo \texttt{.tex}:\\
        \texttt{\textbackslash incluirdiagrama\{Imagens/arch\_visao\_implantacao.png\}\{Visão Física\}\{fig:visao_implantacao_posto\}}%
    \vspace{1cm}}}
    \caption{Visão Física de Implantação e Conectividade Industrial (Placeholder StarUML v7.0)}
    \label{fig:visao_implantacao_posto}
\end{figure}

\section{Catálogo de Registros de Decisão Arquitetural (ADR)}

\subsection{ADR-01: Adoção de Arquitetura Edge-First com SQLite Local}
\paragraph{Status:} Aprovado e Homologado.
\paragraph{Decisores:} Diogo Santos Pires Jandiroba (Lead Architect).
\paragraph{Contexto:} Os postos de abastecimento não podem interromper suas vendas sob nenhuma hipótese decorrente de quedas de link de internet, atrasos na SEFAZ ou latência da nuvem.
\paragraph{Decisão:} Implementar arquitetura \textit{Edge-First}, armazenando todo o estado transacional em banco SQLite local configurado em modo WAL (\textit{Write-Ahead Logging}). Os dados são sincronizados de forma eventual e assíncrona com o PostgreSQL na nuvem via mensageria idempotente.
\paragraph{Consequências Positivas:} Tolerância total a partição de rede; latência sub-milissegundo para autorização de abastecimento na pista.
\paragraph{Consequências Negativas:} Necessidade de resolução de conflitos de conciliação de cadastros centralizados vindos da nuvem.

\subsection{ADR-02: Comunicação com Concentrador via Sockets TCP/IP com Fallback Serial}
\paragraph{Status:} Aprovado.
\paragraph{Contexto:} O ecossistema de postos possui concentradores modernos com placa Ethernet nativa e concentradores legados operando via conversor Serial USB/RS-232.
\paragraph{Decisão:} Projetar o módulo de telemetria utilizando o padrão de projeto \textit{Adapter}, provendo abstração unificada para conexões Socket TCP assíncronas e suporte secundário a portas seriais nativas.
\paragraph{Consequências Positivas:} Flexibilidade de instalação em postos modernos e legados sem alteração no núcleo do sistema.

\section{Governança de Dados, Criptografia e LGPD}

Em estrito atendimento à Lei Geral de Proteção de Dados (Lei Federal nº 13.709/2018):
\begin{itemize}
    \item Os CPFs informados para inclusão no cupom fiscal são anonimizados nos bancos de auditoria analítica;
    \item Os logs de transações financeiras e abastecimentos possuem retenção legal obrigatória de 5 anos;
    \item Toda comunicação entre o Edge Server e a nuvem transita criptografada via protocolo TLS 1.3.
\end{itemize}
